% GENERATED FILE. Edit canonical lessons, then run npm run generate:books.
% canonical-source-hash: fnv1a64:c4ba599e2d2d3bc6
% canonical-lessons: TE-C334-jananam, TE-C334-pasikandu, TE-C334-namakaranam, TE-C334-janmincu, TE-C334-divincu

\chapter{A Baby Is Born}
\label{ch:te-a-baby-is-born}
\hlchaptermodality{334}

\begin{chapteropening}
\textbf{By the end of this chapter:} \emph{I can say a birth, a newborn baby, a naming ceremony, to take birth and to bless.}

Complete the last lesson of chapter 334: I can say a birth, a newborn baby, a naming ceremony, to take birth and to bless.
\end{chapteropening}

\section[\texorpdfstring{\te{జననం} (jananaṁ)}{జననం (jananaṁ)}]{\te{జననం} (jananaṁ) — a birth}

\label{lesson:TE-C334-jananam}



\begin{quote}
Before the new one: say the Telugu for a trace, then the Telugu for a fingerprint.
\end{quote}

\subsection*{You'll want to know: \te{జననం}}
\textbf{\te{జననం}} — \emph{jananaṁ} — \textquotedblleft{}a birth\textquotedblright{}.

\begin{quote}
\te{ఆ} \te{పాప} \te{జననం} \te{ఆసుపత్రిలో} \te{జరిగింది}.
\end{quote}

\begin{quote}
\emph{ā pāpa jananaṁ āsupatrilō jarigindi.}
\end{quote}

\begin{quote}
The baby's birth took place in the hospital.
\end{quote}

A formal word for birth: \te{జనన} \te{తేదీ} is the date of birth on a form.

\begin{cousinweb}
\textbf{\te{జననం}} (\emph{jananaṁ}) is borrowed from Sanskrit \emph{janana}, \textquotedblleft{}birth\textquotedblright{} (CDIAL 5099), from the root \emph{jan}, \textquotedblleft{}to be born\textquotedblright{}.
\end{cousinweb}

\begin{grammarlens}[title={What you've built}]
One more word for telling what happened.
\end{grammarlens}

\subsection*{Your turn}
\begin{itemize}
\raggedright
  \item \emph{Say it:} \emph{jananaṁ}
  \item \emph{Say it:} \emph{jananaṁ}, once more
  \item \emph{Say it:} \emph{ā pāpa jananaṁ āsupatrilō jarigindi.} — The baby's birth took place in the hospital.
  \item \emph{Recall:} say the Telugu for a trace, then the Telugu for a fingerprint, then say \emph{jananaṁ} again
\end{itemize}

\begin{tcolorbox}[breakable,colback=teal!4,colframe=teal!35!black,title={Before you move on}]
What does \te{జననం} mean? (\textquotedblleft{}A birth\textquotedblright{}.) Say it once more.
\end{tcolorbox}
\section[\texorpdfstring{\te{పసికందు} (pasikandu)}{పసికందు (pasikandu)}]{\te{పసికందు} (pasikandu) — a newborn baby}

\label{lesson:TE-C334-pasikandu}



\begin{quote}
Before the new one: say the Telugu for a fingerprint, then the Telugu for a birth.
\end{quote}

\subsection*{You'll want to know: \te{పసికందు}}
\textbf{\te{పసికందు}} — \emph{pasikandu} — \textquotedblleft{}a newborn baby\textquotedblright{}.

\begin{quote}
\te{పసికందు} \te{రాత్రంతా} \te{ఏడ్చింది}.
\end{quote}

\begin{quote}
\emph{pasikandu rātrantā ēḍcindi.}
\end{quote}

\begin{quote}
The newborn cried all night.
\end{quote}

Plural: \te{పసికందులు}, \emph{pasikandulu}.

A tender word for a baby only days or weeks old.

\begin{cousinweb}
\textbf{\te{పసికందు}} (\emph{pasikandu}) is a Telugu compound whose second part, \emph{kandu}, \textquotedblleft{}an infant\textquotedblright{}, is native Dravidian (DEDR 1411, with Tamil \emph{kaṉṟu}, \textquotedblleft{}a calf, a young one\textquotedblright{}).
\end{cousinweb}

\begin{grammarlens}[title={What you've built}]
One more word for telling what happened.
\end{grammarlens}

\subsection*{Your turn}
\begin{itemize}
\raggedright
  \item \emph{Say it:} \emph{pasikandu}
  \item \emph{Say it:} \emph{pasikandu}, once more
  \item \emph{Say it:} \emph{pasikandu rātrantā ēḍcindi.} — The newborn cried all night.
  \item \emph{Recall:} say the Telugu for a fingerprint, then the Telugu for a birth, then say \emph{pasikandu} again
\end{itemize}

\begin{tcolorbox}[breakable,colback=teal!4,colframe=teal!35!black,title={Before you move on}]
What does \te{పసికందు} mean? (\textquotedblleft{}A newborn baby\textquotedblright{}.) Say it once more.
\end{tcolorbox}
\section[\texorpdfstring{\te{నామకరణం} (nāmakaraṇaṁ)}{నామకరణం (nāmakaraṇaṁ)}]{\te{నామకరణం} (nāmakaraṇaṁ) — a naming ceremony}

\label{lesson:TE-C334-namakaranam}



\begin{quote}
Before the new one: say the Telugu for a birth, then the Telugu for a newborn baby.
\end{quote}

\subsection*{You'll want to know: \te{నామకరణం}}
\textbf{\te{నామకరణం}} — \emph{nāmakaraṇaṁ} — \textquotedblleft{}a naming ceremony\textquotedblright{}.

\begin{quote}
\te{పాపకు} \te{ఆదివారం} \te{నామకరణం} \te{చేశారు}.
\end{quote}

\begin{quote}
\emph{pāpaku ādivāraṁ nāmakaraṇaṁ cēśāru.}
\end{quote}

\begin{quote}
They held the baby's naming ceremony on Sunday.
\end{quote}

\te{నామం}, a name, and \te{కరణం}, a making: \te{నామకరణం} \te{చేయు}, to name a child.

\begin{cousinweb}
\textbf{\te{నామకరణం}} (\emph{nāmakaraṇaṁ}) is borrowed from Sanskrit \emph{nāmakaraṇa}, \textquotedblleft{}the ceremony of naming after birth\textquotedblright{} (CDIAL 7066), from \emph{nāman}, \textquotedblleft{}a name\textquotedblright{}, and \emph{karaṇa}, \textquotedblleft{}a making\textquotedblright{}.
\end{cousinweb}

\begin{grammarlens}[title={What you've built}]
One more word for telling what happened.
\end{grammarlens}

\subsection*{Your turn}
\begin{itemize}
\raggedright
  \item \emph{Say it:} \emph{nāmakaraṇaṁ}
  \item \emph{Say it:} \emph{nāmakaraṇaṁ}, once more
  \item \emph{Say it:} \emph{pāpaku ādivāraṁ nāmakaraṇaṁ cēśāru.} — They held the baby's naming ceremony on Sunday.
  \item \emph{Recall:} say the Telugu for a birth, then the Telugu for a newborn baby, then say \emph{nāmakaraṇaṁ} again
\end{itemize}

\begin{tcolorbox}[breakable,colback=teal!4,colframe=teal!35!black,title={Before you move on}]
What does \te{నామకరణం} mean? (\textquotedblleft{}A naming ceremony\textquotedblright{}.) Say it once more.
\end{tcolorbox}
\section[\texorpdfstring{\te{జన్మించు} (janmiñcu)}{జన్మించు (janmiñcu)}]{\te{జన్మించు} (janmiñcu) — to take birth; to be born}

\label{lesson:TE-C334-janmincu}



\begin{quote}
Before the new one: say the Telugu for a newborn baby, then the Telugu for a naming ceremony.
\end{quote}

\subsection*{You'll want to know: \te{జన్మించు}}
\textbf{\te{జన్మించు}} — \emph{janmiñcu} — \textquotedblleft{}to take birth; to be born\textquotedblright{}.

\begin{quote}
\te{నేను} \te{ఒక} \te{చిన్న} \te{పల్లెలో} \te{జన్మించాను}.
\end{quote}

\begin{quote}
\emph{nēnu oka cinna pallelō janmiñcānu.}
\end{quote}

\begin{quote}
I was born in a small village.
\end{quote}

Past (I …): \te{జన్మించాను}, \emph{janmiñcānu}.

A formal partner of \te{పుట్టు}: life stories and the news use \te{జన్మించు}.

\begin{cousinweb}
\textbf{\te{జన్మించు}} (\emph{janmiñcu}) is built on Sanskrit \emph{janman}, \textquotedblleft{}birth\textquotedblright{} (CDIAL 5113), with the Telugu verb ending \emph{-iñcu}.
\end{cousinweb}

\begin{grammarlens}[title={What you've built}]
One more word for telling what happened.
\end{grammarlens}

\subsection*{Your turn}
\begin{itemize}
\raggedright
  \item \emph{Say it:} \emph{janmiñcu}
  \item \emph{Say it:} \emph{janmiñcu}, once more
  \item \emph{Say it:} \emph{nēnu oka cinna pallelō janmiñcānu.} — I was born in a small village.
  \item \emph{Recall:} say the Telugu for a newborn baby, then the Telugu for a naming ceremony, then say \emph{janmiñcu} again
\end{itemize}

\begin{tcolorbox}[breakable,colback=teal!4,colframe=teal!35!black,title={Before you move on}]
What does \te{జన్మించు} mean? (\textquotedblleft{}To take birth; to be born\textquotedblright{}.) Say it once more.
\end{tcolorbox}
\section[\texorpdfstring{\te{దీవించు} (dīviñcu)}{దీవించు (dīviñcu)}]{\te{దీవించు} (dīviñcu) — to bless}

\label{lesson:TE-C334-divincu}



\begin{quote}
Before the new one: say the Telugu for a naming ceremony, then the Telugu for to take birth.
\end{quote}

\subsection*{You'll want to know: \te{దీవించు}}
\textbf{\te{దీవించు}} — \emph{dīviñcu} — \textquotedblleft{}to bless\textquotedblright{}.

\begin{quote}
\te{అమ్మమ్మ} \te{పసికందును} \te{దీవించింది}.
\end{quote}

\begin{quote}
\emph{ammamma pasikandunu dīviñcindi.}
\end{quote}

\begin{quote}
Grandmother blessed the newborn.
\end{quote}

Past (I …): \te{దీవించాను}, \emph{dīviñcānu}.

Its noun is \te{దీవెన}, a blessing: elders bless, the young bow.

\begin{cousinweb}
\textbf{\te{దీవించు}} (\emph{dīviñcu}) is a Telugu verb whose origin is not settled; DEDR has no entry for it.
\end{cousinweb}

\begin{grammarlens}[title={What you've built}]
One more word for telling what happened.
\end{grammarlens}

\subsection*{Your turn}
\begin{itemize}
\raggedright
  \item \emph{Say it:} \emph{dīviñcu}
  \item \emph{Say it:} \emph{dīviñcu}, once more
  \item \emph{Say it:} \emph{ammamma pasikandunu dīviñcindi.} — Grandmother blessed the newborn.
  \item \emph{Recall:} say the Telugu for a naming ceremony, then the Telugu for to take birth, then say \emph{dīviñcu} again
\end{itemize}

\begin{tcolorbox}[breakable,colback=teal!4,colframe=teal!35!black,title={Before you move on}]
What does \te{దీవించు} mean? (\textquotedblleft{}To bless\textquotedblright{}.) Say it once more.
\end{tcolorbox}
